% Unit 031 terjemahan Bahasa Indonesia dari chapter3.tex baris 9--55.
% Sumber: Methods of Algebra, Volume 2, commit
% 9a5803ff77dd3257484cb177f851a73770a59dd3 (CC BY 4.0).
% stable-unit-id: o014.aljabr2.chapter3.overview
% Cakupan: ikhtisar Bab 3 sebelum Bagian 3.1.

% segment-id: o014.li2.chapter3.overview.h001
\chapter{Kompleks}\label{sec:cplx}
\hypertarget{o014.aljabr2.chapter3.overview}{ }

% segment-id: o014.li2.chapter3.overview.p001
Konsep kompleks telah diperkenalkan secara ringkas dalam
\S~\ref{sec:cohomology}. Untuk kompleks $X=(X^n,d^n)_n$ dalam kategori
abelian $\mathcal{A}$, kohomologi
$\Hm^n(X):=\Ker(d^n)/\Image(d^{n-1})$%
\footnote{Catatan penerjemah (O014-C032): sumber mencetak
$\Image(d^{n+1})$. Dengan konvensi kokrantai
$X^{n-1}\xrightarrow{d^{n-1}}X^n\xrightarrow{d^n}X^{n+1}$, subobjek yang
berada di dalam $\Ker(d^n)$ ialah $\Image(d^{n-1})$; target memulihkan indeks
tersebut, sesuai dengan definisi sebelumnya dalam buku ini.}
mengekstrak invarian penting dari $X$. Sebaliknya,%
\footnote{Catatan penerjemah (O014-C030): sumber mencetak tanda titik yang
diikuti satu karakter tersesat sebelum klausa tentang lenyapnya semua
kohomologi, sehingga transisinya tidak gramatikal. Konteksnya menyatakan arah
kebalikan; target memulihkan ``Sebaliknya'' tanpa mengubah klaim matematis.}
kompleks yang seluruh kohomologinya nol merupakan barisan eksak, yang juga
disebut kompleks asiklik.

% segment-id: o014.li2.chapter3.overview.p002
Banyak materi dalam
\S\S~\sourcecrossref{sec:Hom-cplx}{3.2}--%
\sourcecrossref{sec:double-cplx}{3.5} pada bab ini tidak melibatkan
kohomologi dan berlaku bagi kategori aditif umum $\mathcal{A}$. Kategori
kompleks pada $\mathcal{A}$ dinotasikan dengan $\cate{C}(\mathcal{A})$.
Kecuali dinyatakan lain, semua kompleks yang dibahas dalam bab ini adalah
kompleks kokrantai sebagaimana disebutkan dalam
Catatan~\ref{rem:cochain-vs-chain}; setiap kali $\Bbbk$ disebutkan, simbol itu
 menyatakan sembarang gelanggang komutatif.

% segment-id: o014.li2.chapter3.overview.p003
Untuk kompleks $X$ dan $Y$ pada kategori aditif $\mathcal{A}$, homotopi
merupakan relasi ekuivalensi pada
$\Hom_{\cate{C}(\mathcal{A})}(X,Y)$. Relasi ini juga dapat dipahami melalui
kompleks $\Hom$, yakni $\Hom^\bullet(X,Y)$, yang dikaji dalam
\S~\sourcecrossref{sec:Hom-cplx}{3.2}. Jika $\mathcal{A}$ kategori abelian,
morfisme-morfisme yang homotopik menginduksi peta yang sama pada kohomologi.
Selanjutnya, kerucut pemetaan $\Cone(f)$ yang diperkenalkan dalam
\S~\sourcecrossref{sec:mapping-cone}{3.3} akan memainkan peran penting dalam
kajian kategori turunan. Kerucut pemetaan berkaitan erat dengan homotopi, dan
keduanya mempunyai penafsiran topologis langsung.

% segment-id: o014.li2.chapter3.overview.p004
Bikompleks
$X=(X^{p,q},\dhori^{p,q},\dvert^{p,q})_{p,q}$ dapat dibayangkan sebagai
kompleks dengan dua arah yang saling bebas, atau secara ekuivalen sebagai
``kompleks dari kompleks''. Bikompleks $X$ dapat diratakan melalui jumlah
langsung (atau hasil kali langsung) menjadi kompleks $\tot_{\oplus}X$ (atau
$\tot_{\Pi}X$), yang disebut kompleks total dari $X$. Definisi tepatnya
melibatkan sejumlah tanda positif dan negatif; rincian itu merupakan pokok
bahasan \S~\sourcecrossref{sec:double-cplx}{3.5}.

% segment-id: o014.li2.chapter3.overview.p005
Salah satu teknik dasar aljabar homologis ialah menurunkan barisan eksak
panjang dalam kohomologi dari barisan eksak pendek kompleks
$0\to X'\xrightarrow{f}X\xrightarrow{g}X''\to0$:
% segment-id: o014.li2.chapter3.overview.q001
\[
  \cdots \to \Hm^n(X') \xrightarrow{\Hm^n(f)} \Hm^n(X)
  \xrightarrow{\Hm^n(g)} \Hm^n(X'')
  \xrightarrow{\text{morfisme penghubung}} \Hm^{n+1}(X') \to \cdots .
\]
Barisan eksak panjang juga dapat dipahami dari sudut pandang kerucut
pemetaan, tetapi hal ini memerlukan argumen yang tidak sepenuhnya sepele.
Semua itu dibahas dalam
\S\S~\sourcecrossref{sec:Abel-cplx}{3.6}--%
\sourcecrossref{sec:cone-vs-long-exact-sequence}{3.7}.

% segment-id: o014.li2.chapter3.overview.p006
Untuk membantu pembaca membiasakan diri dengan berbagai operasi dasar pada
kompleks, \S~\sourcecrossref{sec:HH}{3.8} menyajikan latihan terpadu melalui
homologi Hochschild $\HHm_n(M)$ dan kohomologi Hochschild $\HHm^n(M)$ dari
$M$, suatu bimodul $(R,R)$. Homologi dan kohomologi Hochschild mempunyai
banyak kegunaan dalam matematika; bab ini hanya memberikan pengantar ringkas
untuk tujuan pembelajaran, dan keduanya akan kembali muncul sebagai contoh
dalam bab-bab berikutnya.

% segment-id: o014.li2.chapter3.overview.p007
Berdasarkan funktor pemenggalan yang diperkenalkan dalam
\S~\sourcecrossref{sec:truncation-functors}{3.9}, dalam
\S~\sourcecrossref{sec:double-cplx-coh}{3.10} kita akan menghubungkan
kohomologi vertikal atau horizontal suatu bikompleks dengan kohomologi
kompleks totalnya. Bagian ini memerlukan syarat keterbatasan tertentu pada
bikompleks. Beberapa buku ajar menangani sifat-sifat tersebut dengan barisan
spektral; buku ini mengikuti \cite{KS06} dan menempuh jalur yang agak memutar
untuk membuktikannya secara langsung.

% segment-id: o014.li2.chapter3.overview.p008
Untuk objek $X$ dalam kategori abelian $\mathcal{A}$, resolusi $X$ mempunyai
versi kiri dan kanan, yang secara konkret merupakan barisan eksak%
\footnote{Catatan penerjemah (O014-C031): sumber mencetak
$0\to X\to I^0\to I^2\to\cdots$, sehingga melewatkan $I^1$. Diagram yang
langsung mengikutinya memakai $I^0,I^1,I^2$ secara berurutan; target
memulihkan suku $I^1$.}
% segment-id: o014.li2.chapter3.overview.q002
\[
  0\to X\to I^0\to I^1\to I^2\to\cdots
  \quad\text{atau}\quad
  \cdots\to P_1\to P_0\to X\to0 .
\]
Hal ini juga berarti mengganti $X$ dengan kompleks yang sifatnya lebih baik
melalui kuasi-isomorfisme $X\to I:=(I^n)_n$ atau $P:=(P_n)_n\to X$, dengan
$P$ suatu kompleks rantai. Kuasi-isomorfisme adalah morfisme yang menginduksi
isomorfisme pada tingkat kohomologi atau homologi. Jika setiap $I^n$ (atau
$P_n$) disyaratkan merupakan objek injektif (atau projektif), konstruksi itu
disebut resolusi injektif (atau projektif) dari $X$. Resolusi akan dikaji
secara mendalam dalam \S~\sourcecrossref{sec:resolutions}{3.11}. Dalam arti
berikut, resolusi juga funktorial dan unik hingga homotopi: untuk versi
injektif, misalkan bagian
garis penuh dari diagram berikut berada dalam $\mathcal{A}$.
% segment-id: o014.li2.chapter3.overview.g001
\[
\begin{tikzcd}
	0 \arrow[r] & X' \arrow[r] & (I')^0 \arrow[r] & (I')^1 \arrow[r]
	  & (I')^2 \arrow[r] & \cdots \\
	0 \arrow[r] & X \arrow[u,"f"] \arrow[r]
	  & I^0 \arrow[dashed,u,"{\beta^0}"] \arrow[r]
	  & I^1 \arrow[r] \arrow[dashed,u,"{\beta^1}"]
	  & I^2 \arrow[r] \arrow[dashed,u,"{\beta^2}"] & \cdots
\end{tikzcd}
\]
Andaikan kedua baris eksak dan setiap $(I')^n$ merupakan objek injektif.
Terdapat anak panah putus-putus $\beta^0,\beta^1,\ldots$ yang membuat seluruh
diagram komutatif, dan morfisme kompleks $\beta:I\to I'$ unik hingga
homotopi. Pernyataan ini merupakan penerapan bersama
Teorema~\sourcecrossref{prop:resolution-homotopy}{pada \S~3.11} dan
Lema~\sourcecrossref{prop:ext-alpha-beta}{pada \S~3.11}, yang merupakan kasus
khusus dari pernyataan yang lebih luas.

% segment-id: o014.li2.chapter3.overview.p009
Sebenarnya, \S~\sourcecrossref{sec:resolutions}{3.11} tidak hanya membahas
resolusi objek $X$ dalam $\mathcal{A}$, tetapi juga resolusi kompleks.
Hasil-hasilnya mencakup resolusi Cartan--Eilenberg yang praktis
(Teorema~\sourcecrossref{prop:CE-resolution}{pada \S~3.11}). Sejumlah argumen
 untuk versi kompleks tak terhindarkan menjadi panjang, meskipun gagasan
dasarnya tetap jelas.

% segment-id: o014.li2.chapter3.overview.p010
Dengan mengganti objek kategori $\mathcal{A}$ oleh resolusi injektif (atau
projektif), menerapkan funktor eksak kiri (atau eksak kanan)
$F:\mathcal{A}\to\mathcal{B}$ pada setiap suku resolusi, lalu mengambil
$\Hm^n$ (atau $\Hm_n=\Hm^{-n}$), kita sampai pada definisi klasik funktor
turunan kanan $\mathrm{R}^nF$ (atau funktor turunan kiri $\mathrm{L}_nF$).
Barisan eksak pendek dalam $\mathcal{A}$ kemudian menginduksi barisan eksak
panjang funktor turunan; lihat
\S~\sourcecrossref{sec:derived-primer}{3.12}. Dengan resolusi kompleks,
funktor turunan juga dapat dihitung pada kompleks terbatas bawah atau terbatas
atas; dalam literatur klasik, konstruksi ini disebut funktor hiperturunan.

% segment-id: o014.li2.chapter3.overview.p011
Dari funktor turunan kanan (atau kiri) beserta barisan eksak panjangnya muncul
secara alami konsep funktor-$\delta$ kohomologis (atau homologis), dan funktor
turunan bersifat ``universal'' di antara funktor-funktor tersebut.
Proposisi~\sourcecrossref{prop:erasable-univ-delta}{pada \S~3.12}
mencirikan funktor-$\delta$ universal melalui syarat dapat dihapus (atau dapat
dihapus secara dual) dalam
Definisi~\sourcecrossref{def:erasable}{pada \S~3.12}. Hasil ini merupakan
sumbangan Grothendieck. Perangkat teknik yang berkaitan dengannya membentuk
inti aljabar homologis klasik.

% segment-id: o014.li2.chapter3.overview.p012
Konstruksi penting lainnya ialah funktor turunan kanan (atau kiri) dari
bifunktor $F:\mathcal{A}_1\times\mathcal{A}_2\to\mathcal{B}$, dengan syarat
$F$ eksak kiri (atau eksak kanan) dalam kedua variabel. Contoh yang lazim
ialah funktor $\Ext$ dan $\Tor$ dalam
\S~\sourcecrossref{sec:Ext-Tor}{3.14}; yang pertama selalu melibatkan kompleks
pada kategori lawan. Teorema~\sourcecrossref{prop:balanced-primer}{pada
\S~3.14} mengkaji kelas bifunktor yang disebut bifunktor seimbang dan
menunjukkan bahwa penurunan terhadap salah satu variabel memberikan hasil
yang sama. Pembuktiannya bertumpu pada sifat-sifat dasar kohomologi
bikompleks.

% segment-id: o014.li2.chapter3.overview.p013
Andaikan $\mathcal{A}$ mempunyai hasil kali terhitung yang eksak.
Konstruksi $\lim^1$ yang dibahas dalam
\S~\sourcecrossref{sec:lim1}{3.13} merupakan contoh lain funktor turunan
kanan: ia adalah funktor turunan kanan pertama dari $\varprojlim$, sedangkan
semua funktor turunan kanan yang lebih tinggi bernilai nol
(Teorema~\sourcecrossref{prop:lim1}{pada \S~3.13}). Untuk memperoleh kriteria
praktis yang menjamin $\lim^1=0$, kita akan memperkenalkan syarat
Mittag--Leffler dalam Definisi~\sourcecrossref{def:ML}{pada \S~3.13}.

% segment-id: o014.li2.chapter3.overview.p014
Terakhir, dalam \S~\sourcecrossref{sec:K-injectives}{3.15} kita akan mengkaji
kompleks K-injektif dan K-projektif, lalu mendefinisikan resolusi K-injektif
dan K-projektif. Keduanya masing-masing merupakan perumuman alami resolusi
injektif dan projektif untuk kompleks tak terbatas. Karena K-injektivitas dan
K-projektivitas merupakan sifat kompleks secara keseluruhan, definisinya lebih
ringkas dan alami daripada resolusi injektif dan projektif yang ditentukan
suku demi suku. Dengan konsep tersebut, definisi funktor turunan dapat
diperluas ke kompleks tak terbatas. Persoalannya ialah keberadaan resolusi
semacam itu: Teorema~\sourcecrossref{prop:K-injective-resolution}{pada
\S~3.15} dan Teorema~\sourcecrossref{prop:K-projective-resolution}{pada
\S~3.15} memberikan sejumlah hasil parsial. Argumennya relatif rumit dan
memerlukan beberapa sifat $\lim^1$.

% segment-id: o014.li2.chapter3.overview.n001
\begin{petunjukbacaan}
	Pengalaman menangani kompleks modul-$\Bbbk$ ($\Bbbk$ suatu gelanggang
	komutatif) membantu pembaca memahami isi bab ini dengan lebih cepat, tetapi
	tidak mutlak diperlukan. Berdasarkan Bab~\ref{sec:Abel-cat}, menangani
	kategori abelian yang lebih umum daripada kategori modul juga tidak
	menimbulkan kesulitan mendasar.

	Pembaca yang ingin segera memahami definisi klasik funktor turunan, tetapi
	tidak berminat pada motivasi topologisnya, disarankan terlebih dahulu melewati
	bagian tentang kerucut dan silinder pemetaan dalam
	\S~\sourcecrossref{sec:mapping-cone}{3.3},
	\S~\sourcecrossref{sec:opposite-cplx}{3.4}, dan
	\S~\sourcecrossref{sec:cone-vs-long-exact-sequence}{3.7}. Selain itu,
	\S~\sourcecrossref{sec:opposite-cplx}{3.4} bertujuan menghubungkan
	$\cate{C}(\mathcal{A}^{\opp})$ dengan
	$\cate{C}(\mathcal{A})^{\opp}$. Sejumlah tanda positif dan negatif di sana
	perlu ditangani dengan cermat, tetapi tidak banyak memengaruhi gagasan
	besar. Pada pembacaan pertama, pembaca disarankan hanya mempelajari
	pernyataannya dan melewati pembuktiannya.

	Karena materi berikutnya tidak bergantung pada teori homologi dan
	kohomologi Hochschild, pembaca dapat memutuskan untuk melewatinya berdasarkan
	waktu dan minat. Materi itu tidak sulit dan akan berulang kali muncul sebagai
	contoh atau catatan dalam bab-bab berikutnya.

	Demikian pula, isi \S~\sourcecrossref{sec:lim1}{3.13} hanya muncul sebagai
	contoh dalam bagian-bagian selanjutnya, kecuali dalam
	\S~\sourcecrossref{sec:K-injectives}{3.15}. Isi
	\S~\sourcecrossref{sec:K-injectives}{3.15} hanya dirujuk dalam
	\S~\sourcecrossref{sec:unbounded-derived}{4.11} dan
	\S~\sourcecrossref{sec:otimesL}{4.12}. Jika tidak terlalu berminat pada
	funktor turunan tak terbatas, pembaca dapat terlebih dahulu mempelajari
	bagian definisi dan Contoh~\sourcecrossref{eg:bdd-K-resolution}{pada
	\S~3.15}, lalu melewati bagian selebihnya.
\end{petunjukbacaan}
